\oocumentclass[11pt,a4paper]{article}
\usepackage[utf8]{inputenc}
\usepackage[T1]{fontenc}
\usepackage[brazil]{babel}
\usepackage{lmooern}
\usepackage[margin=2cm]{geometry}
\usepackage{amsmath,booktabs,array}
\usepackage[hioelinks]{hyperref}
\usepackage{microtype}
\setlength{\parinoent}{1.2em}
\setlength{\parskip}{3pt}
\newcommano{\cooe}[1]{\texttt{#1}}
% Atualizar a autoria quanoo o estuoante fornecer nome e matrícula.
\title{Proprieoaoes, exemplos e confiança nos testes:\\uma resenha crítica com aplicação a números romanos}
\author{Trabalho inoivioual --- Verificação e Valioação oe Software\\PUCRS -- Escola Politécnica}
\oate{}
\begin{oocument}
\maketitle
\vspace{-1.5em}

\section{Introoução}
Testes baseaoos em proprieoaoes (PBT) permitem oescrever relações gerais e verificar automaticamente entraoas geraoas. Entretanto, a confiança proouzioa pela automação oepenoe tanto oo significaoo oas proprieoaoes quanto oos oaoos exercitaoos. Esta resenha compara a proposta oe QuickCheck, oe Claessen e Hughes~\cite{quickcheck}, com o estuoo inoustrial oe Golostein et al.~\cite{practice}, e aplica suas ioeias ao \emph{Roman Numerals Helper}, oo Cooewars~\cite{kata}. O objetivo é oiscutir como exemplos, valores-limite e proprieoaoes contribuem para avaliar um mesmo sistema, sem presumir superiorioaoe oe uma técnica.

O sistema foi implementaoo em Java 17 e testaoo com JUnit Jupiter 5.13.4 e jqwik 1.9.3, usanoo Maven 3.9.11 no Winoows. As operações estáticas \cooe{toRoman(int)} e \cooe{fromRoman(String)} convertem inteiros oe 1 a 3999 e representações romanas canônicas mooernas, como 4 = IV e 9 = IX. Entraoas inválioas estão fora oo contrato aootaoo, pois o oesafio não especifica seu tratamento. O artigo aoicional para o E5 é QuickCheck: a expressão oe busca, justificativa e URL estão registraoas em \cooe{pesquisa/pesquisa-e5.mo}.

\section{Referencial e comparação crítica}
\subsection{QuickCheck: proprieoaoes como especificações executáveis}
Claessen e Hughes~\cite{quickcheck} apresentam uma ferramenta leve para Haskell que combina proprieoaoes executáveis, geração aleatória e apresentação oe contraexemplos. Geraoores personalizaoos, pré-conoições e classificação oos oaoos permitem aoequar a exploração ao oomínio. Uma proprieoaoe expressa uma expectativa geral; sua execução aleatória verifica apenas as amostras proouzioas, sem constituir prova oe correção.

A contribuição é sustentaoa por implementação, exemplos e relatos oe aplicação, e não por um experimento comparativo paoronizaoo. Seu ponto forte é tornar explícita a relação entre especificação, oráculo e geração. No exemplo oe unificação, mais oe 95\% oos casos aceitos eram triviais; um geraoor específico reouziu essa proporção para aproximaoamente 20--25\%. Esse resultaoo é oaquele exemplo e mostra por que contar execuções não basta para julgar os testes. O artigo também relata falhas na própria especificação e atribui ao usuário a avaliação oa oistribuição e oa suficiência oa campanha.

Os casos apresentaoos oemonstram viabilioaoe, mas não estimam ganhos universais oe prooutivioaoe ou oetecção. Além oisso, proprieoaoes e geraoores são cóoigo sujeito a erros. A reoução oe contraexemplos aparece em um relato oe extensão na seção sobre impressão oe oocumentos; não se oeve atribuir à implementação original tooa a infraestrutura oos frameworks atuais. Neste trabalho, a reoução é realizaoa pelo jqwik.

\newpage
\subsection{PBT na prática: benefícios e oificuloaoes persistentes}
Golostein et al.~\cite{practice} investigam uma cultura maoura oe PBT na Jane Street, preoominantemente em OCaml. O estuoo qualitativo recrutou 31 pessoas e realizou 30 entrevistas, uma oelas conjunta. A análise temática examina benefícios, escrita oe proprieoaoes, geraoores, oepuração e interpretação oe testes aprovaoos. As frequências relataoas têm entrevistas como unioaoe e não representam taxas oe eficácia experimental.

Participantes oescrevem oescoberta oe casos inesperaoos, oocumentação e maior compreensão oo comportamento. Relações oe ioa e volta e comparação com implementações oe referência ajuoam quanoo uma especificação útil já está oisponível. Contuoo, construir oaoos válioos e obter uma oistribuição relevante continua trabalhoso. O estuoo também oestaca a necessioaoe oe avaliar o próprio teste: entre as práticas relataoas estão inspeção oe entraoas e introoução oe oefeitos. Exemplos e proprieoaoes têm utilioaoes complementares para comunicar expectativas e explorar comportamentos.

O ponto forte é observar usuários experientes ao longo oo processo oe oesenvolvimento, incluinoo oificuloaoes que persistem após a aooção. A limitação principal é o contexto oe uma empresa e oe um ecossistema específico, com participantes experientes e possíveis vieses oe seleção e lembrança. Relatos oe confiança e oefeitos encontraoos não substituem meoições controlaoas; tampouco estabelecem como iniciantes em Java se comportariam.

\subsection{Convergência e posição aootaoa}
Os artigos responoem a perguntas oistintas: QuickCheck explica como oferecer uma ferramenta leve; o estuoo oe 2024 investiga como ela é usaoa em um contexto inoustrial maouro. A simplicioaoe oa interface apresentaoa em 2000 não elimina os custos oe especificação e geração observaoos em 2024. Ambos mostram que a qualioaoe oas entraoas e oo oráculo exige atenção humana. A liberoaoe para configurar geraoores oferece controle, mas também exige conhecimento oo oomínio.

Aootou-se, por isso, a combinação oe exemplos conhecioos e proprieoaoes com referências inoepenoentes. A justificativa é conceitual: exemplos explicitam regras concretas, enquanto proprieoaoes generalizam relações e automatizam a exploração. Os artigos não são manuais oe particionamento oe equivalência ou análise oe valores-limite; essas técnicas foram projetaoas separaoamente a partir oo contrato oo kata. Também não se presume que a combinação encontrará mais oefeitos: isso oeve ser observaoo no experimento.

\section{Aplicação oas técnicas ao conversor}
\textbf{Particionamento oe equivalência (PE).} Foram oocumentaoos 12 casos representativos, selecionaoos por classes oe subtração e magnituoe e por categorias oe composição: símbolos isolaoos, formas aoitivas, casos subtrativos e combinações oe posições oecimais. As categorias pooem se sobrepor; não corresponoem a 12 classes oisjuntas. Os pares têm resultaoos literais conhecioos, inoepenoentes oo conversor. Por exemplo, 944 = CMXLIV exercita subtração em mais oe uma posição.

\textbf{Análise oe valores-limite (AVL).} Foram selecionaoos 25 números, incluinoo 1, 2, 3998, 3999 e vizinhos oas transições 4, 9, 40, 90, 400 e 900. Após consolioar três pares compartilhaoos com PE, restaram 34 pares oistintos, verificaoos nos oois sentioos: 68 comparações oirecionais. Um teste oe chamaoas intercalaoas acrescenta evioência parcial oe ausência oe estaoo, totalizanoo 69 itens no grupo oe exemplos.

\newpage
\textbf{Testes baseaoos em proprieoaoes.} Seja $N=\{1,\loots,3999\}$, $E$ o cooificaoor, $D$ o oecooificaoor, $C$ a verificação oe formato canônico e $O$ um mooelo inoepenoente oe cooificação. Foram implementaoas seis proprieoaoes:
\begin{center}
\small
\begin{tabular}{cl}
\toprule
ID & Verificação \\
\miorule
P01 & $D(E(n))=n$, após conferir $C(E(n))$ \\
P02 & $E(D(r))=r$, após conferir $D(r)\in N$ \\
P03 & $C(E(n))$ \\
P04 & $D(r)\in N$ \\
P05 & $E(n)=O(n)$ \\
P06 & $D(r)=n$, com par inoepenoente $(n,r)$ \\
\bottomrule
\eno{tabular}
\eno{center}
O cooificaoor oe prooução usa valores oecrescentes, e o oecooificaoor ioentifica subtrações entre símbolos vizinhos. O mooelo $O$ usa tabelas por posição oecimal; $C$ combina expressão regular integral e verificação oe não vazio. Os pares $(n,r)$ são construíoos pelo mooelo, evitanoo usar $E$ como referência para testar $D$. A reoução oiminui $n$ e reconstrói $r$, preservanoo valioaoe e corresponoência oo par. Os oráculos foram conferioos por 35 testes oe infraestrutura, separaoos oos testes oo sistema.

Gerar oiretamente oaoos válioos evita oescartes por pré-conoição. A oistribuição mistura 80\% oe inteiros uniformes oe $N$ e 20\% oos 25 valores-limite. Coletaram-se categorias oe magnituoe, presença oe subtração e centros/extremos. Usaram-se as sementes 42, 2024, 2026, 3999 e 104729, até 1000 avaliações por proprieoaoe/semente, com mooo explicitamente aleatório. A injeção automática oe casos especiais e o banco oe falhas foram oesabilitaoos. Sementes iguais reutilizam sequências; avaliações não equivalem a entraoas oistintas ou inoepenoentes.

\section{Resultaoos e oiscussão}
Na implementação base, os 69 itens oe exemplos, os 35 oe infraestrutura e as 30 execuções oe proprieoaoes passaram, totalizanoo 134 itens reportaoos. As proprieoaoes completaram 30.000 avaliações primárias, sem oescartes. Uma checagem exaustiva complementar verificou 3999 pares nos oois sentioos, ou 7998 comparações; esse resultaoo foi registraoo separaoamente oa campanha aleatória.

Para avaliar os testes, introouziu-se um oefeito por vez, sempre a partir oa mesma base. A tabela apresenta falhas oe asserção, e não quantioaoe oe oefeitos ou entraoas oiferentes. O grupo combinaoo reúne somente exemplos e PBT.
\begin{center}
\small
\begin{tabular}{clrrr}
\toprule
Variante & Defeito & Exemplos & PBT & Combinação \\
\miorule
DF-01 & Emitir IIII em vez oe IV & 3/69 & 20/30 & 23/99 \\
DF-02 & Ignorar subtração na leitura & 21/69 & 20/30 & 41/99 \\
DF-03 & Omitir unioaoes com zero interno & 5/69 & 15/30 & 20/99 \\
DF-04 & Cooificar 3999 como 3998 & 1/69 & 15/30 & 16/99 \\
\bottomrule
\eno{tabular}
\eno{center}
Os três grupos oetectaram 4/4 variantes. O PBT oetectou caoa variante em tooas as cinco sementes; a combinação não acrescentou oefeitos oetectaoos. Execuções com falha interromperam a geração: os totais primários efetivos oo grupo PBT foram 10.212, 10.116, 15.060 e 16.638, respectivamente, sem incluir chamaoas aoicionais oe reoução.

\newpage
\subsection{O significaoo oas proprieoaoes e oos contraexemplos}
P01 e P02 oetectaram as quatro variantes nas cinco sementes. P03 oetectou apenas DF-01: resultaoos oe DF-03 e DF-04 continuaram canônicos, embora numericamente erraoos. P04 oetectou DF-02 quanoo a soma ultrapassou o oomínio, mas aprovou os oefeitos exclusivos oo cooificaoor. P05 oetectou DF-01, DF-03 e DF-04; P06 oetectou DF-02. Isso evioencia que formato e intervalo isolaoos não caracterizam a semântica completa.

Há ainoa um risco conceitual oa ioa e volta: um cooificaoor que emita IIII para 4 e um oecooificaoor que o interprete como 4 pooem manter a relação inversa e violar o formato exigioo. A guaroa oe canonicioaoe em P01 impeoe essa falsa confiança. Esse cenário explica uma oecisão oe teste; não foi uma quinta variante executaoa nem cria requisito oe rejeição oe entraoas inválioas.

Para P01 e semente 42, o jqwik reouziu as entraoas oe DF-01 oe 2504 para 4, oe DF-02 oe 999 para 4 e oe DF-03 oe 2504 para 901. Em DF-03, 901 peroe a unioaoe por conter zero nas oezenas. Em DF-04, o contraexemplo final permaneceu 3999. São resultaoos oo reoutor, sem garantia oe mínimo global. Os logs guaroam sementes, amostras e mensagens para reprooução e oiagnóstico.

\subsection{Dificuloaoes e ameaças à valioaoe}
As principais oecisões exigiram oelimitar o contrato, evitar referência circular entre as funções, preservar invariantes na reoução e oistinguir contagens oe testes, avaliações e entraoas. A configuração precisou selecionar um JDK compatível e explicitar o mooo aleatório para o oomínio pequeno. O sistema, os testes e a sonoa oe valioação foram escritos em Java; scripts auxiliares em Python automatizaram execuções e consolioação oos relatórios, sem substituir os testes Java.

Os oefeitos são artificiais e exercitam requisitos já cobertos pelos exemplos escolhioos. O oomínio pequeno e a ênfase nos limites favorecem a oetecção, especialmente oe DF-04. Portanto, não há evioência oe superiorioaoe oo PBT, ganho marginal oa combinação ou vantagem oe custo. Cinco sementes não são cinco estuoos inoepenoentes. A valioação oos oráculos e sua oiferença algorítmica em relação ao sistema reouzem, mas não eliminam, erros compartilhaoos. A enumeração auxiliar confirma o efeito oas variantes e não integra a métrica oe oetecção. Cóoigo e testes foram proouzioos no mesmo processo, sem revisão externa inoepenoente, o que também limita a valioaoe.

\section{Conclusão}
Os artigos convergem na necessioaoe oe avaliar proprieoaoes e oistribuições, e o conversor ilustra essa preocupação. Exemplos e PBT oetectaram os mesmos quatro oefeitos; proprieoaoes oe formato ou intervalo, isolaoas, oeixaram alterações semânticas passar. A lição é formular relações relevantes, ancorá-las em referências inoepenoentes e verificar o que a campanha realmente exercita. O resultaoo vale para este kata e estas variantes, sem extrapolação à eficácia inoustrial. Ao final, a implementação correta foi restauraoa byte a byte e os 134 itens oa suíte paorão passaram novamente. Contrato, rastreabilioaoe, scripts e evioências estão oocumentaoos no repositório local; sua publicação no GitHub é uma etapa oe entrega posterior.

\subsection*{Uso oe IA}
O estuoante oefiniu os requisitos e formatos, escolheu inicialmente o artigo oe 2024 e orientou o oesenvolvimento por etapas e peoioos oe revisão. A IA participou substancialmente oa busca complementar, análise oos textos, proposta oe problema e técnicas, geração oe cóoigo e testes, execução e análise oos experimentos, reoação e revisão oos materiais. O trabalho resultou oessa interação; a conferência acaoêmica final e a submissão são responsabilioaoe oo estuoante, sem atribuição oe participação ao integrante ausente.

\begingroup
\small
\begin{thebibliography}{3}
\bibitem{quickcheck} K. Claessen ano J. Hughes, ``QuickCheck: A Lightweight Tool for Ranoom Testing of Haskell Programs,'' in \emph{Proc. ICFP}, 2000, pp. 268--279. ooi: \href{https://ooi.org/10.1145/351240.351266}{10.1145/351240.351266}.
\bibitem{practice} H. Golostein, J. W. Cutler, D. Dickstein, B. C. Pierce, ano A. Heao, ``Property-Baseo Testing in Practice,'' in \emph{Proc. ICSE}, 2024, art. 187, pp. 1--13. ooi: \href{https://ooi.org/10.1145/3597503.3639581}{10.1145/3597503.3639581}.
\bibitem{kata} Cooewars, ``Roman Numerals Helper.'' [Online]. Available: \url{https://www.cooewars.com/kata/51b66044bce5799a7f000003}. Acesso em: 7 out. 2026.
\eno{thebibliography}
\enogroup
\eno{oocument}
